\section{整除、最大公约数与 Euclid 算法}
\label{sec:04-Discrete-Mathematics/04-Number-Theory/01}

给你两个整数 252 和 105，怎么找到它们的最大公约数？你可以把两个数分别分解质因数：\(252=2^2\cdot3^2\cdot7\)，\(105=3\cdot5\cdot7\)，公共部分是 \(3\cdot7=21\)。但如果换成 1234567 和 890123，分解质因数就远没有这么顺手了。可 Euclid 在 2300 年前找到的办法，根本不依赖质因数分解——它只反复做带余除法，余数越除越小，最后一个非零余数就是答案。

这个算法的力量来自一个简单事实：整除关系的运作方式，让公因数在``取余数''这一步纹丝不动。在展开算法之前，先把这个事实讲清楚。

\subsection{整除是一条存在整数因子的关系}

\(a\mid b\) 的定义：存在整数 \(k\) 使 \(b=ak\)。重点不是``\(b/a\) 除得尽''这个模糊感觉，而是那个见证 \(k\) 的具体存在。比如 \(3\mid 15\)，见证是 \(k=5\)；\(7\mid 0\)，见证是 \(k=0\)（任何数都整除 0）。有了 \(k\) 这个整数，它就能参与后续的代数运算——你可以把它跟别的见证加加减减，推导出新的整除关系。

这个定义方式的价值在下面这条性质中立刻体现。若 \(a\mid b\) 且 \(a\mid c\)，则对任意整数 \(x,y\)，
\[
a\mid xb+yc.
\]

证明只需写两行：设 \(b=ak_1\)、\(c=ak_2\)，则
\[
xb+yc = a(xk_1+yk_2),
\]
新见证是 \(xk_1+yk_2\)，它仍是一个整数。整除对整数线性组合封闭——别小看这一行代入，它是整个 Euclid 算法和 Bézout 恒等式下面那块最底层的砖。

拿具体数字确认思路：\(7\mid 21\) 且 \(7\mid 35\)，那你取 \(4\cdot21 - 2\cdot35 = 14\)，7 照样整除 14。无论怎么拿 21 和 35 的整数倍拼凑，7 永远整除结果。因为 7 同时是这两个数的因子，它自然会留在它们的任何整数线性组合里。

反过来想想什么情况下这条性质不能用来传递整除信息。如果只知 \(6\mid 18\) 而不知 6 是否整除另一个数，那 \(18x + \text{?}y\) 是否被 6 整除就取决于问号的值了。整除对线性组合的封闭性需要一个数同时整除组合中的每一项——这是接下来 Euclid 算法中每次约化都依赖的前提。

\subsection{最大公约数是所有公因数的代表}

\(d=\gcd(a,b)\) 是同时整除 \(a,b\) 的最大正整数。若 \(\gcd(a,b)=1\)，称 \(a,b\) 互素。注意互素不要求二者本身是素数——8 与 15 各自不是素数，但最大公约数是 1，它们互素。

现在回到开头的困境：两个不小的整数，怎么高效地求最大公约数？质因数分解的路面对大数时走不通——目前还没有已知的快速分解整数的算法。换个角度想一想：如果把一个大数减掉另一个数的若干倍，公因数会丢失吗？

先看一个草稿版的思路，不用除法，只用减法。252 和 105：从 252 中减掉 105，得到 147；再减 105，得到 42。公因数在每一步减法中都保留着——因为 \(147=252-105\)，整除 252 和 105 的因子也整除 147。做到底就是反复从较大的数中减去较小的数，直到两个数相等。这就是最朴素的 Euclid 算法形式：反复做减法，公因数纹丝不动。

但减法版有个效率问题。如果两个数相差很远——比如 1000000 和 7——你需要连减十几万次。带余除法一步跳到余数：\(1000000=7\cdot142857+1\)，直接得到 1。公因数保留的逻辑和减法一致，只不过一次跨了更多步。

回到 252 和 105。105 和 42 的公因数有哪些？3 整除 105（\(105=3\cdot35\)）和 42（\(42=3\cdot14\)），7 也整除二者。和 252 与 105 的公因数一样。再约化一次：105 和 42，\(105=42\cdot2+21\)，余数 21。42 和 21——21 整除 42（\(42=21\cdot2\)），最大公约数是 21。和分解质因数得到的答案一致。这个草稿揭示的不是巧合。

把上面的草稿提炼成精确形式，就是除法算法：
\[
a=bq+r,\qquad 0\le r<|b|.
\]
\(a\) 除以 \(b\)，商 \(q\)，余数 \(r\) 严格小于 \(|b|\)。比如 \(252\div105\)：\(q=2\)，\(r=42\)。

现在到了最关键的一步。\(a\) 和 \(b\) 的公因数集合，是否恰好等于 \(b\) 和 \(r\) 的公因数集合？验证两个方向：若 \(d\mid a\) 且 \(d\mid b\)，由 \(r=a-bq\) 和整除对线性组合的封闭性，\(d\mid r\)。反过来，若 \(d\mid b\) 且 \(d\mid r\)，由 \(a=bq+r\)，同样得 \(d\mid a\)。公因数集合完全相同，所以：
\[
\gcd(a,b)=\gcd(b,r).
\]

这一步把求两个大数的 gcd 变成了求两个更小数的 gcd。除数变被除数，余数变除数，问题规模持续缩小：
\[
\begin{aligned}
252 &= 105\cdot 2 + 42,\\[2pt]
105 &= 42\cdot 2 + 21,\\[2pt]
42 &= 21\cdot 2 + 0.
\end{aligned}
\]

余数序列 \(252\to105\to42\to21\to0\) 严格递减，每一步都完整保留了公因数集合。最后一个非零余数 21 就是 \(\gcd(252,105)\)。\(\gcd(a,0)=a\)（任何非零整数整除 0），所以当余数归零时，当时的除数就是答案。

为什么算法一定终止？余数序列是非负整数，且每一步严格减小。自然数不能无限递减——这是自然数最小性（良序原理）的直接推论。有限步内余数必然归零。这个终止保证不依赖任何超出整数基本性质的东西。

\subsection{回代得到 Bézout 系数}

Euclid 算法向前走，找到了 gcd。往回走呢？从倒数第二步的等式开始：
\[
21 = 105 - 2\cdot 42.
\]
但第二步告诉我们 \(42 = 252 - 2\cdot 105\)，代入：
\[
21 = 105 - 2(252 - 2\cdot105) = 5\cdot 105 - 2\cdot 252.
\]
\(\gcd(252,105)=21\) 被写成了 \((-2)\cdot252 + 5\cdot105\)。系数 \(-2\) 和 \(5\) 就是一对 Bézout 系数。

一般地，对任意整数 \(a,b\)，存在整数 \(x,y\) 使
\[
\gcd(a,b)=ax+by.
\]

这个恒等式也可以换个方向理解，不依赖回代过程。考虑所有形如 \(ax+by\)（\(x,y\) 遍历整数）的整数构成的集合。因为 gcd 整除 \(a\) 和 \(b\)，它一定整除这个集合中的每一个数——换句话说，集合中的元素全都是 gcd 的倍数。所以这个集合里最小的正整数至少是 gcd。而回代过程构造出的 \(x,y\) 恰好取到了 gcd 本身。于是 gcd 正好等于所有整数线性组合中的最小正整数。Bézout 恒等式不是额外安上去的一条性质——它是 gcd 的另一种等价描述。

这个视角顺手回答了另一个问题。线性 Diophantine 方程
\[
ax+by=c
\]
何时有整数解？若 \(\gcd(a,b)\mid c\)，取 Bézout 等式 \(ax_0+by_0=\gcd(a,b)\)，两边同乘 \(c/\gcd(a,b)\) 即得一组解。反之，若存在整数解，则 gcd 整除左边 \(ax+by\)，因而整除右边 \(c\)。条件充要：\(ax+by=c\) 有整数解当且仅当 \(\gcd(a,b)\mid c\)。

扩展 Euclid 算法是回代过程的系统化版——在逐次取余的同时一路维护系数，不需要先走到底再回头代入。每步维护两个辅助序列：记录当前余数如何由最初的 \(a,b\) 线性表示。算法结束时，最后一组系数就是 Bézout 系数。复杂度随输入位数增长很慢（与输入位数的平方成正比），是模逆计算、RSA 密钥生成等密码操作的底层工具。

找出一组 Bézout 系数后，方程 \(ax+by=\gcd(a,b)\) 的全部整数解随之确定。若 \((x_0,y_0)\) 是一组特解，则通解为
\[
x = x_0 + \frac{b}{\gcd(a,b)}t,\qquad y = y_0 - \frac{a}{\gcd(a,b)}t,\quad t\in\Z.
\]
验证：代入 \(ax+by\)，含 \(t\) 的项恰好抵消。比如对 \(252x+105y=21\)，特解 \(x=-2,y=5\)，通解为 \(x=-2+5t,\;y=5-12t\)。任取 \(t=1\)：\(x=3,y=-7\)，验算 \(252\cdot3+105\cdot(-7)=756-735=21\)。这就是线性 Diophantine 方程的通解结构，来源就是 Bézout 恒等式。

一个工程上的小提醒：手工回代可行，但写代码时应该直接用扩展 Euclid 算法。回代需要先记下所有中间商，再用递归或迭代的方式一层层代回去；扩展 Euclid 在线性时间中同时维护 gcd 和系数，避免了存储开销。

\subsection{Euclid 引理与唯一分解}

素数有一个特别的``穿透''性质：若素数 \(p\mid ab\)，则 \(p\mid a\) 或 \(p\mid b\)。这个引理值得停下来想一想——它并不平凡。比如 \(6\mid 2\cdot3\)，但 6 既不整除 2 也不整除 3，因为 6 是合数。素数之所以有这个性质，恰好是因为它``没有多余因子''：一个素数与你所给的数之间，要么整除，要么互素，没有第三种可能。

证明用 Bézout 恒等式。若 \(p\nmid a\)，则 \(\gcd(p,a)=1\)。为什么？\(p\) 是素数，它的正因数只有 1 和 \(p\)。若 \(\gcd(p,a) > 1\)，则 \(\gcd(p,a)=p\)，与 \(p\nmid a\) 矛盾。于是存在整数 \(x,y\) 使 \(px+ay=1\)。两边同乘 \(b\)：
\[
pxb + aby = b.
\]
\(pxb\) 被 \(p\) 整除（含有因子 \(p\)）；\(aby\) 因已知条件 \(p\mid ab\) 也被 \(p\) 整除。左侧两项都被 \(p\) 整除，因此右侧 \(b\) 也被 \(p\) 整除。所以 \(p\mid b\)。

注意这个证明里 Bézout 恒等式做了最关键的一步：把``互素''翻译成``存在线性组合等于 1''。只有这种翻译之后，\(p\mid ab\) 的条件才能通过代数操作传到 \(b\) 身上。没有 Bézout 这一步翻译，素数对乘积的穿透力无法从 gcd 的定义直接推出。

这条引理是算术基本定理的支柱。每个大于 1 的整数可以唯一分解为素数的乘积（不计顺序）。分解的存在性用强归纳就够了：要么 \(n\) 本身是素数，要么 \(n=ab\)，而 \(a,b<n\) 由归纳假设皆可分解。唯一性则需要 Euclid 引理来保证：假设一个整数有两种分解式，取第一种分解中的任一素数 \(p\)，它整除整个乘积；反复应用 Euclid 引理，迫使 \(p\) 必须整除第二种分解中的某个素数因子；而两个素数若互相整除，它们必然相等。于是两边的素数集合完全相同，计数也一致。

我们从``两个数怎么找最大公约数''这个具体问题出发，一路走到了整数算术的骨架。Euclid 算法不只是找 gcd 的快捷方式——向前走，找到 gcd；向后走，把 gcd 变成两个数的整数线性组合。这个组合表示反过来证明了素数对乘积的穿透力，而素数的穿透力又支撑了整个整数的唯一分解定理。一条链上的三个环节，缺了哪一环后面的都推不下去。
